\section{Choose an additional observation deliberately}
The recipe ending in \texttt{-1} is each authored step's primary observation.
The additional recipes below keep that same cumulative program and vary its inputs
or caller controls. Prepare each in a new source file. A changed input can change
every working value, so do not compare a marker-state result with the hexadecimal
values printed for the RFC example.

A marker seed starts with state words zero through fifteen. At the load checkpoint,
those values reveal row and lane positions directly. After column arithmetic they
have already changed; a later shuffle moves those changed values. The diagram labels
\emph{state positions}, not numerical values that remain zero through fifteen
throughout the calculation. Compare a selected lane with the corresponding lane
before selection, or with the independently calculated recipe observation.

\begin{longtable}{lp{0.28\linewidth}p{0.47\linewidth}}
  \caption{Additional exact recipes and their observation purpose.}                                              \\
  \toprule Recipe & Prepared variation                   & Observation                                           \\\midrule\endfirsthead
  \toprule Recipe & Prepared variation                   & Observation                                           \\\midrule\endhead
  \texttt{s04-2}  & Four all-one words                   & Word carry is discarded in the first group.           \\
  \texttt{s07-2}  & Four all-one words                   & Check all four final working words before stores.     \\
  \texttt{s24-2}  & Marker state                         & Loaded rows retain the seed's index order.            \\
  \texttt{s27-2}  & Marker state                         & Selected rows use the correct column-round positions. \\
  \texttt{s30-2}  & Marker state                         & Ten pairs end in restored state order.                \\
  \texttt{s33-2}  & Zero bytes                           & Return zero without supplied pointer accesses.        \\
  \texttt{s33-3}  & 4097 bytes                           & Return one before a frame or output store.            \\
  \texttt{s33-4}  & \(2^{32}\) bytes                     & Check the full count; return one.                     \\
  \texttt{s34-2}  & 65 bytes, last counter               & Return two before output stores.                      \\
  \texttt{s41-2}  & One byte                             & Transform only one destination byte.                  \\
  \texttt{s41-3}  & 64 bytes, last counter               & The last available block succeeds.                    \\
  \texttt{s41-4}  & 65 bytes, last counter               & A required second block is rejected.                  \\
  \texttt{s41-5}  & 128 bytes, counter \texttt{fffffffe} & Use exactly the last two blocks.                      \\
  \texttt{s41-6}  & 129 bytes, counter \texttt{fffffffe} & Reject the required third block.                      \\
  \texttt{s41-7}  & 4096 bytes in place                  & Check full-block repetition at the local size limit.  \\
  \texttt{s41-8}  & 65 bytes in place                    & Read each original byte before replacing it.          \\
  \texttt{s41-9}  & Zero bytes                           & Return zero without allocating a frame.               \\
  \texttt{s41-10} & 4097 bytes                           & Reject the count and preserve caller memory.          \\
  \bottomrule
\end{longtable}

Here the last counter means hexadecimal \texttt{ffffffff}. Requests that return
early have no internal arithmetic checkpoint to visit. Pause at \texttt{main\_done}
in the caller to inspect their status before it is stored; the function view pauses
at its supplied return continuation. The input files contain only the actual bounded
fixture bytes. A declared rejected count such as \(2^{32}\) does not allocate or
claim to process a message of that size.
